\dictchar{ई}

\bhojentry{ई}{\K{\bii}}{Ee}{स्वर / सर्वनाम | तद्भव}{\starsThree}%
{हिंदी तथा संस्कृत का चौथा स्वर वर्ण; यह, ई।}%
{ई (Vowel EE / Pronoun: यह)}%
{Fourth vowel of Indic alphabet; demonstrative pronoun 'this'.}
\bhojexample{ई हमार गाँव बा।}{This is my village.}

\bhojentry{ईमान}{\K{\bii\bma\bmaa\bna}}{Eemaan}{हि. पुं. | अरबी आगत}{\starsThree}%
{धर्म, निष्ठा, नीयत, सच्चाई।}%
{ईमान / निष्ठा}%
{Faith, integrity, honesty.}

\bhojsubentry{ईमानदार}{\K{\bii\bma\bmaa\bna\bda\bmaa\bra}}{Eemaandaar}{हि. वि. | अरबी-फारसी}{\starsThree}%
{सच्चा, निष्ठावान व्यक्ति।}%
{ईमानदार}%
{Honest, trustworthy person.}

\bhojsubentry{ईमानदारी}{\K{\bii\bma\bmaa\bna\bda\bmaa\bra\bmii}}{Eemaandaari}{हि. स्त्री. | अरबी-फारसी}{\starsThree}%
{सच्चाई, निष्ठा, सत्यनिष्ठा।}%
{ईमानदारी}%
{Honesty, integrity.}

\bhojsubentry{ईमान-धर्म}{\K{\bii\bma\bmaa\bna}-\K{\bdha\bra\bma}}{Eemaan-dharm}{हि. पुं. | अरबी-संस्कृत}{\starsThree}%
{नीत और धर्म।}%
{ईमान-धर्म}%
{Moral conscience and faith.}

\bhojentry{ईंट}{\K{\bii\banus\btta}}{Eent}{हि. स्त्री. | तद्भव}{\starsThree}%
{भट्ठे में पकाई मिट्टी की चौकोर सिल्ल।}%
{ईंट (Brick)}%
{Brick.}
\bhojexample{ईंट से नया मकान बनल बा।}{A new house has been built with bricks.}

\bhojsubentry{ईंट-भट्ठा}{\K{\bii\banus\btta}-\K{\bbha\bvir\btta\btha\bmaa}}{Eent-bhattha}{हि. पुं. | तद्भव}{\starsThree}%
{ईंट पकाने की भट्ठी।}%
{ईंट के भट्ठा}%
{Brick kiln.}

\bhojsubentry{ईंट-चोट}{\K{\bii\banus\btta}-\K{\bca\bmo\btta}}{Eent-chot}{हि. स्त्री. | [ठेठ]}{\starsTwo}%
{ईंट-पत्थर फेंक कर मारपीट करना।}%
{ईंट-चोट / पत्थरबाजी}%
{Stone pelting, brickbatting.}

\bhojentry{ईख}{\K{\bii\bkha}}{Eekh}{हि. स्त्री. | तद्भव}{\starsThree}%
{गन्ना, ऊख।}%
{ईख (Sugarcane) / ऊख}%
{Sugarcane.}
\bhojsee{ऊख}

\bhojsubentry{ईख के रस}{\K{\bii\bkha} \ \K{\bka\bme} \ \K{\bra\bsa}}{Eekh ke ras}{हि. पुं. | [ठेठ]}{\starsThree}%
{गन्ने का मीठा रस।}%
{ईख के रस}%
{Sugarcane juice.}

\bhojentry{ईसर}{\K{\bii\bsa\bra}}{Eesar}{हि. पुं. | तद्भव}{\starsThree}%
{ईश्वर, परमात्मा, भगवान।}%
{ईसर / ईश्वर / भगवान}%
{God, Supreme Being.}

\bhojsubentry{ईसरी}{\K{\bii\bsa\bra\bmii}}{Eesari}{हि. स्त्री. | [ठेठ]}{\starsTwo}%
{देवी, परमेश्वरी शक्ति।}%
{ईसरी / देवी}%
{Goddess, divine power.}

\bhojentry{ईस}{\K{\bii\bsa}}{Ees}{हि. स्त्री. | [ठेठ/कृषि]}{\starsTwo}%
{बैलगाड़ी या हल की मुख्य लंबी पतली काठ।}%
{ईस (Ees) / हल के काठ}%
{Long wooden beam of plough/cart.}

\bhojentry{ईहा}{\K{\bii\bha\bmaa}}{Eeha}{सं. स्त्री. | तत्सम}{\starsTwo}%
{इच्छा, अभिलाषा, चाह; चेष्टा।}%
{इच्छा / चाह}%
{Desire, wish, effort.}

\bhojsubentry{ईहापोह}{\K{\bii\bha\bmaa\bpa\bmo\bha}}{Eehapoh}{सं. पुं. | तत्सम}{\starsTwo}%
{उधेड़बुन, असमंजस, सोच-विचार।}%
{उधेड़बुन / असमंजस}%
{Perplexity, inner deliberation.}

\bhojentry{ईर्ष्या}{\K{\bii\bra\bvir\bssa\bvir\bya\bmaa}}{Eershya}{सं. स्त्री. | तत्सम}{\starsTwo}%
{जलन, डाह, कुढ़न।}%
{डाह (Daah) / जलन}%
{Envy, jealousy.}

\bhojsubentry{ईर्ष्यालु}{\K{\bii\bra\bvir\bssa\bvir\bya\bmaa\bla\bmu}}{Eershyaalu}{सं. वि. | तत्सम}{\starsTwo}%
{डाहू, जलने वाला व्यक्ति।}%
{डाहू / जलैहा}%
{Envious person.}

\bhojentry{ईद}{\K{\bii\bda}}{Eed}{हि. स्त्री. | अरबी आगत}{\starsThree}%
{मुसलमानों का पवित्र खुशी का त्योहार।}%
{ईद (Eid)}%
{Eid festival.}

\bhojsubentry{ईदगाह}{\K{\bii\bda\bga\bmaa\bha}}{Eedgaah}{हि. पुं. | अरबी-फारसी}{\starsThree}%
{ईद की सामूहिक नमाज़ पढ़ने का खुला मैदान।}%
{ईदगाह}%
{Eid prayer ground.}

\bhojsubentry{ईदी}{\K{\bii\bda\bmii}}{Eedi}{हि. स्त्री. | अरबी आगत}{\starsTwo}%
{ईद पर बड़ों द्वारा छोटों को दिया जाने वाला उपहार या पैसा।}%
{ईदी}%
{Eid gift/money.}

\bhojentry{ईस्वी}{\K{\bii\bsa\bvir\bva\bmii}}{Eesvi}{हि. वि. | अरबी आगत}{\starsThree}%
{ईसा मसीह के जन्म से शुरू ईसवी संवत।}%
{ईसवी (AD / CE)}%
{Christian Era (AD/CE).}

\bhojentry{ईश}{\K{\bii\bsha}}{Eesh}{सं. पुं. | तत्सम}{\starsTwo}%
{स्वामी, राजा, भगवान।}%
{स्वामी / मालिक}%
{Lord, master, ruler.}

\bhojsubentry{ईशान}{\K{\bii\bsha\bmaa\bna}}{Eeshaan}{सं. पुं. | तत्सम}{\starsTwo}%
{उत्तर-पूर्व दिशा का कोण (ईशान कोण)।}%
{उत्तर-पूरब कोना}%
{North-East direction.}

\bhojentry{ईश्वर}{\K{\bii\bsha\bvir\bva\bra}}{Eeshwar}{सं. पुं. | तत्सम}{\starsThree}%
{परमात्मा, सर्वशक्तिमान प्रभु।}%
{ईश्वर / भगवान}%
{God, Almighty Lord.}

\bhojsubentry{ईश्वरीय}{\K{\bii\bsha\bvir\bva\bra\bmi\bya}}{Eeshwariye}{सं. वि. | तत्सम}{\starsTwo}%
{दैवीय, प्रभु का विधान।}%
{दैवीय}%
{Divine, heavenly.}

\bhojentry{ईसा}{\K{\bii\bsa\bmaa}}{Eesa}{हि. पुं. | अरबी आगत}{\starsThree}%
{ईसा मसीह, जीसस क्राइस्ट।}%
{ईसा मसीह}%
{Jesus Christ.}

\bhojsubentry{ईसाई}{\K{\bii\bsa\bmaa\bai}}{Eesaai}{हि. पुं./वि. | अरबी आगत}{\starsThree}%
{ईसाई धर्म के अनुयायी।}%
{ईसाई}%
{Christian.}

\bhojentry{ईंधन}{\K{\bii\banus\bdha\bna}}{Eendhan}{हि. पुं. | तद्भव}{\starsThree}%
{जलावन, बालन, चूल्हा जलाने की सामग्री।}%
{बालन / जलावन}%
{Fuel, firewood.}

\bhojentry{ईही}{\K{\bii\bha\bmii}}{Eehi}{हि. सर्व. | [ठेठ]}{\starsThree}%
{यही (नजदीक की वस्तु या व्यक्ति)।}%
{ईही (Eehi) / यही}%
{This right here, this very thing.}

\bhojentry{ईहो}{\K{\bii\bha\bmo}}{Eeho}{हि. सर्व. | [ठेठ]}{\starsThree}%
{यह भी।}%
{ईहो (Eeho) / यह भी}%
{This as well.}

\bhojentry{ई-सभ}{{\K{\bii}}-\K{\bsa\bbha}}{Ee-sabh}{हि. सर्व. | [ठेठ]}{\starsThree}%
{ये सब (बहुवचन)।}%
{ई-सभ / ये सब}%
{All these.}

\bhojentry{ई-ठवाँ}{{\K{\bii}}-\K{\bttha\bva\banus}}{Ee-thavan}{हि. अव्य. | [ठेठ]}{\starsThree}%
{इस स्थान पर, इसी जगह।}%
{ई-ठवाँ / एह जगह}%
{In this exact place.}
